\clearpage
\setcounter{page}{1}
\maketitlesupplementary
\setcounter{section}{0}
\renewcommand{\thesection}{\Alph{section}}
\renewcommand{\theHsection}{supp.\Alph{section}}

% Section letters A-F are cited by the submitted main paper ("Supp. Sec. A-E"):
% keep their order and content. References to the main paper use \mainref and
% \maineqref (preamble.tex), which print the stored number when this file is
% compiled without the main paper.

This supplement accompanies the main paper. It details the dataset and reference construction (Sec.~\ref{sec:s-data}), the representations and optimization (Sec.~\ref{sec:s-model}), the audit design and estimators with their derivations (Sec.~\ref{sec:s-audit}), the evaluation protocol and additional results (Sec.~\ref{sec:s-eval}), the calibration reference analysis (Sec.~\ref{sec:s-ref}) and a simulation study of the estimator (Sec.~\ref{sec:s-sim}). Equations and sections of the main paper are cited with their numbers there, and Tab.~\ref{tab:notation} collects its notation.

\begin{table}[t]
\centering\small
\setlength{\tabcolsep}{4pt}
\caption{Notation of the main paper. The constants below the rule are fixed design choices (main paper, Sec.~\mainref{sec:audit}{3.2}).}
\label{tab:notation}
\begin{tabular}{@{}ll@{}}
\toprule
Symbol & Meaning \\
\midrule
$X_i$, $Y_i$ & features and reference label of candidate $i$ \\
$m_i$, $p_i$ & weak-mark count, image prediction of a mark \\
$s_i$, $r_i$, $t_0$ & pilot score, within-scan rank, pilot threshold \\
$h_i$, $\bar h_i$ & alert leverage, Eq.~\maineqref{eq:leverage}{1}, floored leverage \\
$u_i$, $v_i$ & working probability and variance, Eq.~\maineqref{eq:working}{2} \\
$\pi_i$, $R_i$ & inclusion probability, Eq.~\maineqref{eq:allocation}{3}, audit indicator \\
$b$, $\mathcal B$ & budget factor (\HalfBudget{} or 1), expected budget \\
$a_j$, $\mathcal I_j$ & reach of target $j$, Eq.~\maineqref{eq:reach}{4}, its candidates \\
$q_i$, $z_i$ & auxiliary, corrected target, Eq.~\maineqref{eq:augmentation}{5} \\
$g$, $S_j$ & linear head, patient score (maximum logit) \\
$\alpha$, $k$, $T$ & alert rate, rank, threshold, Eq.~\maineqref{eq:cutoff}{6} \\
$P$, $U$, $E$ & C/D-negative pool, official-negative subset, $P\setminus U$ \\
\midrule
$\tau$, $\delta$, $\omega$ & leverage constants \RankScale, \ScoreOffset, \ScoreWidth \\
$\beta_m$ & intercepts \WpZeroBase, \WpOneBase, \WpTwoBase{} for $m=0,1,2$ marks \\
$\gamma_m$ & slopes \WpZeroSlope, \WpOneSlope, \WpTwoSlope{} for $m=0,1,2$ marks \\
$\varepsilon_v$, $\varepsilon_h$ & variance floor \VarFloor, leverage floor \LevFloor \\
$\pi_{\min}$, $\pi_{\max}$ & probability floor \ProbFloor${}\times b$, cap \ProbCap \\
$\rho_m$ & mark-stratified probabilities \MarkPiZero, \MarkPiOne, \MarkPiTwo{} ($\times b$) \\
\bottomrule
\end{tabular}
\end{table}

\section{Dataset and Reference Construction}
\label{sec:s-data}

\paragraph{Cohort.}
We indexed all \NScans{} LUNA16 series and verified each MetaImage header and raw volume against the CRC checksum and byte count of the corresponding member of the release archive. Series Instance UIDs were linked to TCIA patient identifiers for the full inventory. The official subsets contain 89 scans each for subsets 0--7 and 88 each for subsets 8 and 9. The TCIA identifiers reveal one patient who appears in both subset 2 and subset 6, of whom we keep the earlier series, and two further scans in subset 2 have eight and three reading sessions, respectively, and are excluded. The remaining 885 series correspond to distinct patients (Tab.~\ref{tab:cohort}). On the test set, the candidate roster covers 274 of the \NTargetTest{} targets, all \NOmitTest{} observed omissions and 150 of the \NOmitBookmark{} omissions imposed by the bookmark cap. Every official nodule annotation is linked to at least one of the four released reading sessions under our geometric matching.

\begin{table}[t]
\centering\small
\caption{Patient splits under the original C/D reference. Neg.\ counts reference-negative patients; Obs.\ omit.\ counts first-reader omissions.}
\label{tab:cohort}
\begin{tabular}{lrrrr}
\toprule
Split & Patients & Neg. & Targets & Obs. omit. \\
\midrule
Train & 445 & 152 & 541 & 92 \\
Cal & 175 & 67 & 215 & 34 \\
Test & 265 & 105 & 276 & 42 \\
\bottomrule
\end{tabular}

\end{table}

\paragraph{Candidate roster.}
The V2 false-positive reduction roster extends the original merged candidate set with the outputs of ETROCAD and M5LCADThreshold0.3~\cite{setio2017,luna16v2}. Its components were developed on several sources, including LIDC-IDRI and, for one of the original systems (LargeCAD), the LUNA16 source data, and some complete CAD submissions used LUNA16 cross-validation; the development exposure of individual V2 coordinates is not available. We discard the supplied class column and use identical coordinates for all methods. The upstream proposals therefore condition both target coverage and the negative tail of the score distribution. Note also that the official FROC evaluation uses a different target definition and ignores some non-target marks, whereas all our counts follow the reference defined in the main paper.

\paragraph{Target geometry.}
For each released contour, we use the positive-inclusion regions, convert pixel coordinates and slice positions to millimetres, and take the midpoint of the bounding box as the centre of the mark. Its size is the largest axis-aligned in-plane extent across slices. Two sets of marks are matched by a one-to-one assignment that minimizes the normalized distance, subject to a maximum distance of $\max(\MinRadiusMM\,\mathrm{mm},(d_1+d_2)/4)$, while maximizing the number of matches. Each matched C/D pair defines a target whose centre and diameter are the means of the two marks. A candidate is assigned to the closest target within one normalized radius, using the scale $\max(\MinRadiusMM\,\mathrm{mm},d/2)$. The omission status of a target is determined by a separate one-to-one matching against the marks of the first weak reader. The signed spacing and origin in the MetaImage header define the correspondence between voxel and world coordinates. The roles and matching rules are identical in all experiments. After requiring four reading sessions, no weak read is missing in this cohort, although the discriminative auxiliary still encodes availability explicitly.

\paragraph{Disputed candidates.}
We call an unmatched candidate \emph{disputed} if it lies near an isolated contour drawn by only one of the two reference readers. The released LIDC reader records are the result of an unblinded second phase~\cite{armato2011}, so an isolated contour indicates disagreement under our operational reference rather than a confirmed finding, and no independent diagnostic adjudication of these candidates was performed. A separate patient-level check found no official positive nodule in subsets 6--9 that had not been exposed under its definition of exposure. All weights and heads were developed without access to scores on subsets 2 and 6--9, and we therefore describe the evaluation as a locked retrospective analysis.

\paragraph{Data use.}
All images and annotations come from the public, de-identified LIDC-IDRI and LUNA16 releases~\cite{armato2011,setio2017}. This is a secondary analysis of these releases: it involved no interaction with participants, no new data collection and no consent procedure, and no new ethics review was conducted for this secondary use.

\FloatBarrier
\section{Representations and Optimization}
\label{sec:s-model}

\paragraph{Hand-crafted features.}
Around each candidate, we interpolate a $13^3$ grid with 1.5\,mm spacing, clip intensities to $[-1100,400]$\,HU and divide by 1000. Four radial shells contribute the mean, the standard deviation and the fraction of voxels above $-300$\,HU, and Gaussian and Hessian filters at four scales contribute a smoothed intensity, the three Hessian eigenvalues and a curvature-based roundness ratio~\cite{li2003}, giving \NDescriptors{} descriptors in total.

\paragraph{Patch encoder.}
The 3D network receives a $\PatchMM^3$ crop resampled to an isotropic 1\,mm grid, with intensities clipped to $[-1050,400]$\,HU, stored with 8-bit precision and scaled to $[-1,1]$. It consists of two convolutional blocks at each of the widths 16, 32 and 48, each with $3^3$ kernels, batch normalization and SiLU activations, and a factor-two pooling after the first two widths. The output concatenates global average pooling with average pooling over the central $4^3$ region, followed by a linear projection to \EmbedDim{} channels, layer normalization and SiLU. During training, a linear readout predicts the weak marks. At inference, we average the embeddings of the original patch and of its flips along each of the two in-plane axes.

The encoder is trained on \NDevPatients{} patients from subsets 0, 1 and 3, using only the marks of the first weak reader, without pretrained weights and without any reference label. Each epoch comprises 72 minibatches of 192 candidates, 32 marked and 160 unmarked. After epochs 4, 10 and 16, we refresh a pool of hard negatives from the top 2.5\% of unmarked scores, from which 48 of the unmarked samples in each batch are drawn, while the remainder are drawn from all unmarked candidates. We use reflections, axis permutations, affine intensity jitter, label smoothing to $0.01/0.99$ and AdamW with a cosine schedule from $1.5{\times}10^{-3}$ to $1.5{\times}10^{-4}$. The epoch-16 checkpoint of a 24-epoch development schedule is frozen, and the linear heads receive its standardized \EmbedDim{}-dimensional embedding. The bookmark-cap setting (Sec.~\ref{sec:s-eval}) uses a separately trained encoder, since the cap removes first-reader marks during encoder training.

\paragraph{Model selection.}
We first evaluated a multi-scale 2.5D model with twelve orthogonal slices at two fields of view, over three saved epochs and two ridge values. It improved the separation of candidates but did not meet our development criterion for all-target localization. We then compared rank-weighted heads with one 3D model per setting, requiring an improvement of ten points in all-target recall and an improvement in within-scan target separation, with negative alerts within five points. The unweighted 3D model satisfied these criteria and was retained.

\paragraph{Linear heads.}
Features are standardized on the fitting candidates, and each head is a linear logistic model with a ridge penalty of $10^{-4}$ on the coefficients and $10^{-6}$ on the intercept. Since the logistic loss is linear in its target, optimizing the augmented loss uses the exact first moment of the corrected labels, and the curvature of the loss with respect to the score does not depend on the target, even when the corrected target is negative; the ridge penalty keeps the problem bounded. IPW applies the logistic loss to audited candidates with weights $1/\pi_i$, normalized by their sum. The nnPU baseline treats the marks of reader A as positives, uses the Horvitz--Thompson estimate of the class prior clipped to the interior of $(0,1)$, and minimizes the convex nonnegative risk. The EM heads use a two-reader working likelihood: starting from the IPW image head, they alternate 16 updates with audited labels held fixed, model the probability of a mark given a reference positive with a logistic head (ridge 0.02) for each reader, and use a smoothed constant for the mark rate given a reference negative.

\FloatBarrier
\section{Audit Design and Estimation}
\label{sec:s-audit}

\paragraph{Weak evidence.}
Within each of three folds of training patients, an image-based logistic model fitted on the other folds predicts whether either weak reader marked a candidate, giving the probability $p_i$ used in Eq.~\maineqref{eq:working}{2}. Its coefficients down-weight concordant marks and emphasize disagreement, and they were fixed before any reference label was revealed. Combining the within-scan rank with an absolute threshold in Eq.~\maineqref{eq:leverage}{1} retains candidates that can determine an extreme scan score even when the intensity distributions of different scans vary.

\paragraph{Allocation rule.}
With $w_i=\bar h_i^2v_i$, the allocation of Sec.~\mainref{sec:audit}{3.2} solves the separable convex problem
\[
 \min_{\pi}\ \sum_i w_i\Big(\frac{1}{\pi_i}-1\Big)
 \quad\text{s.t.}\quad \sum_i\pi_i=\mathcal B,\quad \pi_{\min}b\le\pi_i\le\pi_{\max}.
\]
For a multiplier $\mu>0$ of the budget constraint, each term $w_i/\pi_i+\mu\pi_i$ is convex in $\pi_i$ with its stationary point at $\sqrt{w_i/\mu}$, so its minimizer over the box is $\operatorname{clip}(\sqrt{w_i/\mu},\,\pi_{\min}b,\,\pi_{\max})$. Writing $c=\mu^{-1/2}$ and $\sqrt{w_i}=\bar h_i\sqrt{v_i}$ gives Eq.~\maineqref{eq:allocation}{3}. Since the total $\sum_i\pi_i$ is continuous and nondecreasing in $c$, bisection finds the $c$ that attains the budget whenever $N\pi_{\min}b\le\mathcal B\le N\pi_{\max}$, and by convexity this solution is optimal.

\paragraph{Unbiasedness of the training loss.}
Let $\mathcal F_{-f}$ contain the data available before auditing together with the audit draws of all patient folds other than $f$. Conditional on $\mathcal F_{-f}$ and on the labels $Y$ of the finite pool, the independent draws in fold $f$ satisfy $\Pr(R_i=1)=\pi_i$, while the auxiliary $q_i$ is fixed. The corrected target of Eq.~\maineqref{eq:augmentation}{5} therefore has $\mathbb E_R[z_i]=q_i+\pi_i(Y_i-q_i)/\pi_i=Y_i$ and $\mathrm{Var}_R(z_i)=(1-\pi_i)(Y_i-q_i)^2/\pi_i$. The logistic loss $\ell(z,g)=\log(1+e^{g})-zg$ is affine in $z$. Hence, for any fixed head $g$ and any loss weights $\eta_i$ determined before auditing,
\[
 \mathbb E_R\!\left[\frac{\sum_i \eta_i\ell(z_i,g_i)}{\sum_i \eta_i}\right]
 =\frac{\sum_i \eta_i\ell(Y_i,g_i)}{\sum_i \eta_i}.
\]
This covers both the unit-weight objective ($\eta_i=1$) and the tail-weighted objective ($\eta_i=1+9h_i$), which gives additional weight to candidates near the frozen tail and is compared with the unit-weight objective in Tab.~\ref{tab:contrasts}. The property holds for a misspecified auxiliary, and the weighted estimand does not depend on the inclusion probabilities.

\paragraph{Power tuning.}
Let $\bar X_i$ denote the standardized design vector including the intercept. We scale the auxiliary by a factor $\lambda$ chosen to minimize an estimate of
\[
 \sum_i\frac{1-\pi_i}{\pi_i}\,(Y_i-\lambda q_i)^2\,\|\bar X_i\|^2 .
\] For each outer fold, the auxiliary is predicted on the remaining folds through an inner split, and the criterion is estimated from the audited labels with an additional factor $1/\pi_i$; the minimizer of this quadratic is clipped to $[0,1]$. The held-out fold then uses the scaled prediction of an auxiliary fitted on all other folds. This fixed-fold, scalar variance tuning follows prediction-powered training~\cite{ppssl2025} and uses the same labels and head as the untuned estimator.

\paragraph{Reach and gradient variance.}
The probability $a_j$ that target $j$ is audited follows Eq.~\maineqref{eq:reach}{4}, with an empty product equal to one, and summing $a_j$ over targets or over omissions gives the expected reach. This computation uses the reference labels of the training pool for evaluation only. We additionally record the squared residuals, the extremes of the audited signed residuals and the diagonal design-variance contributions of the corrected label moment. Conditional on the auxiliary fit of a fold, its residual terms are independent under the design; the reported sum of conditional diagonal terms ignores covariances between folds that arise from fitting the auxiliaries on shared labels. We also measure the realized gradient error across draws and the Brier error of each auxiliary by weak-mark pattern (Tab.~\ref{tab:aux}).

\begin{table}[t]
\centering\small
\caption{Auxiliary error by weak-mark pattern. Mean squared residual ($\times 10^3$) of each auxiliary on test candidates at the anchor draw with mark-stratified audits; Marks: weak marks of readers A and B (1: marked); $N$: test candidates; Ref.: reference-positive rate (\%). Computed for evaluation only.}
\label{tab:aux}
\begin{tabular}{lrrrrr}
\toprule
Marks & $N$ & Ref.\% & Image & Disc. & EM \\
\midrule
00 & 224,631 & 0.016 & 0.36 & 0.16 & 0.17 \\
01 & 203 & 27.586 & 200.85 & 190.49 & 550.87 \\
10 & 169 & 14.793 & 128.08 & 116.97 & 573.15 \\
11 & 472 & 70.763 & 253.43 & 197.39 & 292.23 \\
\bottomrule
\end{tabular}

\end{table}

\FloatBarrier
\section{Evaluation Protocol and Additional Results}
\label{sec:s-eval}

\begin{table}[t]
\centering\small
\setlength{\tabcolsep}{4pt}
\caption{Supervision and representation at the anchor draw. Cells give detections among \NOmitTest{} observed omissions / \NTargetTest{} targets / alerts among \PTest{} C/D-negative test patients. All audited methods use the same base-budget mark-stratified draw.}
\label{tab:observed}
\begin{tabular}{lcc}
\toprule
Supervision & Descriptors & 3D features \\
\midrule
First read & 2 / 39 / 12 & 4 / 81 / 2 \\
Second read & 2 / 37 / 12 & 1 / 76 / 3 \\
nnPU & 2 / 40 / 12 & 3 / 59 / 2 \\
Audit-only IPW & 2 / 57 / 14 & 3 / 70 / 5 \\
Paired EM & 2 / 44 / 12 & 3 / 76 / 3 \\
Image augmentation & 2 / 55 / 12 & 4 / 69 / 5 \\
Discriminative augmentation & 2 / 47 / 11 & 1 / 68 / 4 \\
EM augmentation & 2 / 48 / 10 & 2 / 71 / 6 \\
Power-tuned augmentation & 2 / 48 / 11 & 1 / 69 / 4 \\
Full reference & 2 / 47 / 12 & 1 / 74 / 2 \\
\bottomrule
\end{tabular}

\end{table}

\paragraph{Prespecified comparisons.}
The primary half-budget comparison was selected using subset 4 for calibration and subset 5 for evaluation. On these development subsets, the selected method achieved a six-draw omission recall of 23.3\%, compared with 18.3\% for base-budget mark-stratified IPW, both at a mean alert rate of 8.6\%. nnPU with an audited class prior reached 28.3\% over its three development draws and was retained as a separate baseline. A claim of superiority required gains of 10 points over first-read training and 5 points over the strongest simple audited baseline, positive lower confidence limits, and a positive effect against the audited baseline in at least four of the six draws. A claim of efficiency required a saving of at least 25\% of the audited candidates and lower confidence limits of at least $\NIMargin$ points for both recall measures. Both claims further required a realized negative alert rate of at most 13\%, at most 3 points above the audited baseline, and at most 0.20 additional FP/scan. We report the selected comparison irrespective of its outcome.

\paragraph{Supervision and representation.}
Table~\ref{tab:observed} reports every training method with both feature families at the anchor draw, all using the same base-budget mark-stratified audits. The paired effects of the main comparisons and of the auxiliary, weighting and representation choices are in Tab.~\ref{tab:contrasts} and Fig.~\ref{fig:locked}.

\paragraph{Confidence intervals.}
Paired intervals resample test patients \NBoot{} times with fixed weights across the six matched audit draws. Intervals for a single draw condition on that draw and on the pooled calibration set. To assess calibration variability, we resample the \PCal{} negative calibration patients 300 times with replacement and recompute each strict threshold, using the same test patients throughout. Resampling the calibration patients jointly with the audit draws (six draws by 300 resamples), with the test patients held fixed, gives a central \CILevel\% range of $\JointLo$ to $+\JointHi$ points for the omission effect of the selected method relative to base-budget mark IPW (Fig.~\ref{fig:variation}a). The matched-alert analysis uses the negative patients of subset 2 to obtain a grid of thresholds and those of subset 6 to select the highest attained alert rate not exceeding \AlphaPct\%, preferring the higher threshold in case of ties. This analysis is separate from the pooled conformal guarantee.

\begin{figure*}[t]
\centering
\includegraphics[width=\textwidth]{locked_performance}
\caption{Locked operating points and paired effects. (a)~Every point uses a calibration-derived threshold; open markers denote nominal $\alpha=\AlphaNom$. (b)~Paired-tail omission recall minus each comparator with \CILevel\% patient-bootstrap intervals (Tab.~\ref{tab:contrasts}). Filled: nominal threshold; open: calibration-only matched-alert check. The shaded region lies beyond the $\NIMargin$-point margin.}
\label{fig:locked}
\end{figure*}

\begin{table*}[t]
\centering\small
\caption{Paired effects at the locked and calibration-matched operating points. Differences are percentage points except FP/scan; the omission column gives patient-bootstrap \CILevel\% intervals over the matched-draw mean.}
\label{tab:contrasts}
\begin{tabular}{llrrrr}
\toprule
Contrast & Threshold & $\Delta$ Omit. [95\%] & $\Delta$ All & $\Delta$ Alert & $\Delta$ FP/scan \\
\midrule
Selected vs. base IPW & Nominal & -4.0 [-7.0, -1.4] & +3.02 & -1.75 & -0.01 \\
Selected vs. first read & Nominal & -6.3 [-15.1, +0.0] & -6.22 & +0.63 & -0.06 \\
Selected vs. nnPU & Nominal & -4.0 [-11.5, +1.1] & +0.48 & +0.63 & -0.05 \\
Paired vs. mark (same fit, half) & Nominal & -4.4 [-8.3, -0.8] & -1.33 & -1.90 & -0.01 \\
Paired vs. image tail (same fit, half) & Nominal & -1.6 [-3.7, +0.0] & -0.12 & +0.16 & -0.02 \\
Disc. vs. image aug. (mark) & Nominal & -1.2 [-3.0, +0.5] & +0.66 & -0.95 & -0.03 \\
EM vs. disc. aug. (mark) & Nominal & +0.0 [-2.5, +2.1] & -0.36 & +0.32 & -0.00 \\
EM aug. vs. IPW (mark) & Nominal & -0.8 [-5.7, +3.9] & -5.31 & -0.95 & -0.10 \\
EM aug. vs. first read (mark) & Nominal & -4.0 [-13.8, +4.3] & -11.96 & +2.22 & -0.14 \\
EM aug. vs. image aug. (mark) & Nominal & -2.4 [-7.3, +2.3] & -5.68 & -1.27 & -0.10 \\
Power vs. disc. aug. (mark) & Nominal & +0.0 [+0.0, +0.0] & -0.91 & +0.00 & +0.01 \\
Tail-weighted vs. unit-weight disc. & Nominal & -0.4 [-1.9, +1.0] & +2.90 & +0.16 & -0.01 \\
3D vs. descriptors (full ref.) & Nominal & -2.4 [-11.4, +5.6] & +9.78 & -9.52 & +0.02 \\
3D vs. descriptors (first) & Nominal & +4.8 [+0.0, +12.5] & +15.22 & -9.52 & +0.07 \\
Selected vs. base IPW & Matched & -3.6 [-7.0, -0.9] & +0.30 & -1.11 & -0.01 \\
Selected vs. first read & Matched & -7.1 [-17.1, +0.0] & -13.35 & -0.16 & -0.12 \\
Selected vs. nnPU & Matched & -4.8 [-12.7, +0.5] & -5.80 & -0.16 & -0.10 \\
Paired vs. mark (same fit, half) & Matched & -2.4 [-5.7, +0.8] & -2.42 & -1.59 & -0.02 \\
\bottomrule
\end{tabular}

\end{table*}

\begin{figure*}[t]
\centering
\includegraphics[width=\textwidth]{variability}
\caption{Audit, calibration and test-patient variation answer different questions. Paired tail minus mark IPW for (a)~omission recall and (b)~negative-patient alerts. Dark intervals resample test patients with the fitted heads and thresholds fixed, for the anchor fit and for the six-fit mean. Blue points vary the training audit draw with the calibration pool fixed. Green intervals give the central \CILevel\% of calibration resamples at the anchor, and of calibration and audit resampling combined. All use the same test patients.}
\label{fig:variation}
\end{figure*}

\begin{table*}[t]
\centering\small
\caption{Audit allocation and workload, averaged over the matched draws. Candidates: audited candidates; Pos.\ cand.: audited reference-positive candidates; Targets and Omit.: distinct audited targets and omissions; Patients: patients with at least one audit; Exp.\ omit.\ and Exp.\ unmarked: expected reach (Eq.~\maineqref{eq:reach}{4}) over all \NOmitTrain{} training omissions and over the \NOmitUnmarked{} omissions without any weak mark.}
\label{tab:costs}
\begin{tabular}{lrrrrrrrr}
\toprule
Policy & Budget & Candidates & Pos. cand. & Lesions & Omit. & Patients & Exp. omit. & Exp. dbl. unmarked \\
\midrule
Mark stratified & 0.5 & 3,942 & 84 & 75.7 & 14.0 & 443.5 & 13.0 & 0.4 \\
Mark stratified & 1 & 7,933 & 169 & 142.0 & 24.2 & 445.0 & 23.8 & 0.9 \\
Uniform & 0.5 & 3,938 & 9 & 8.5 & 1.0 & 443.3 & 1.4 & 0.5 \\
Uniform & 1 & 7,938 & 15 & 15.0 & 2.3 & 444.8 & 2.7 & 0.9 \\
Image tail & 0.5 & 3,939 & 571 & 431.3 & 65.5 & 445.0 & 64.2 & 22.6 \\
Image tail & 1 & 7,892 & 665 & 490.0 & 80.8 & 445.0 & 79.7 & 29.1 \\
Paired tail & 0.5 & 3,940 & 652 & 482.7 & 79.2 & 444.8 & 79.4 & 27.4 \\
Paired tail & 1 & 7,898 & 689 & 502.2 & 86.2 & 445.0 & 84.6 & 30.2 \\
\bottomrule
\end{tabular}

\end{table*}

\paragraph{Sparse strata and gradient variance.}
At the anchor draw, mark stratification audits none of the 61 reference-positive training candidates without any weak mark, and every audited label in this stratum is multiplied by $1/\pi=50$. With the image-only auxiliary, the realized signed residual correction in this stratum ranges from $-36.51$ to approximately zero, compared with $-0.36$ to approximately zero for the discriminative auxiliary. Using the weak marks thus reduces the auxiliary error in this stratum (Tab.~\ref{tab:aux}), although the deployed head remains limited by both ranking and calibration.

For candidates without any weak mark, the diagonal design variance of the discriminative moment is $5.67{\times}10^{-6}$ under base-budget mark allocation and $1.44{\times}10^{-6}$ under half-budget paired-tail allocation. Weighting the gradient by the frozen leverage $h_i$ gives $3.22{\times}10^{-6}$ and $4.10{\times}10^{-9}$, respectively; this diagnostic is computed on the same tail that guided the allocation. Candidates away from the tail can still receive large inverse probabilities: under half-budget paired-tail allocation, the scaled residuals in this stratum range from $-0.91$ to $8.45$. The power-tuned multipliers at the anchor draw are 0.67, 0.90 and 0.38, and over the six base-budget mark draws the power-tuned and untuned discriminative heads yield identical omission and alert counts at the primary threshold. All-target recall, in contrast, is sensitive to the head objective across the development grid of ridge penalties and loss weights.

\paragraph{Controlled bookmark cap.}
In this setting, the first reader retains only its largest mark in each scan, which creates 151 imposed omissions on the test set. Across the three repeated draws, EM augmentation recovers 29--37 of these omissions, compared with 25 for first-read training, 32--34 for IPW and 27--35 for discriminative augmentation. The corresponding alert counts are 4--5 of 105 negative patients for EM augmentation, 4 for first-read training and 4--6 for both IPW and discriminative augmentation. Since the cap imposes a size-based selection of omissions, it poses a different problem from the observed omissions studied in the main paper.

\paragraph{Per-target changes.}
Averaged over the six draws, the selected method newly detects 9.1\% [7.3, 11.1] of all targets relative to base-budget mark IPW and loses 6.1\% [4.8, 7.5], with paired patient-bootstrap intervals; for the omissions, the corresponding fractions are 0.8\% [0.0, 2.1] and 4.8\% [2.4, 7.8]. At the anchor draw, these changes amount to eight gained and nine lost detections over all targets, including no gained and two lost omissions. In the calibration-matched analysis, the mean test alert rates are 1.7\% for the selected method and 2.9\% for mark IPW. When audit draws and calibration patients are resampled jointly (six draws by 300 resamples), the central 95\% of the omission effect spans $-14.3$ to $+7.1$ points, with the test patients held fixed. All comparisons in this section evaluate different thresholds on the same held-out roster, and no independent adjudication establishes whether the disputed contours correspond to in-scope nodules.


\begin{table}[t]
\centering\small
\setlength{\tabcolsep}{4pt}
\caption{Controlled bookmark cap. Cells give detections at the anchor draw among the \NOmitBookmark{} imposed omissions / \NTargetTest{} targets / alerts among \PTest{} reference-negative patients. The encoder in this setting is trained on the capped first-read labels.}
\label{tab:cap}
\begin{tabular}{lcc}
\toprule
Supervision & Descriptors & 3D features \\
\midrule
First read & 31 / 61 / 16 & 25 / 81 / 4 \\
Second read & 22 / 37 / 12 & 30 / 88 / 5 \\
nnPU & 32 / 62 / 16 & 25 / 81 / 4 \\
Audit-only IPW & 29 / 57 / 15 & 32 / 96 / 4 \\
Paired EM & 24 / 43 / 11 & 33 / 98 / 4 \\
Image augmentation & 28 / 55 / 13 & 33 / 97 / 4 \\
Discriminative augmentation & 26 / 51 / 12 & 34 / 96 / 4 \\
EM augmentation & 24 / 46 / 10 & 37 / 98 / 5 \\
Power-tuned augmentation & 26 / 51 / 13 & 34 / 96 / 4 \\
Full reference & 26 / 47 / 12 & 38 / 104 / 7 \\
\bottomrule
\end{tabular}

\end{table}

\paragraph{Audit workload.}
At half budget, paired-tail sampling audits 3,940 candidates on average, of which 652 are reference positive, whereas base-budget mark stratification audits 7,933 candidates, of which 169 are positive (Tab.~\ref{tab:costs}). Both policies open almost every one of the \NFit{} training patients, so the \AuditSave\% saving concerns candidate decisions; if opening a scan has a fixed cost, the saving in reading time is smaller. Both designs also reuse the existing second read, and the cost of acquiring it, of running the pilot detector and of reviewing a candidate were not measured.

\paragraph{Sources of the acquisition gain.}
The acquisition controls of Sec.~\mainref{sec:reach}{4.2} are exploratory: they were evaluated after the locked study, on the saved half-budget allocation of the training roster, whose reference had already been inspected. The image prediction $p_i$ is recovered from the saved working probability $u_i$ by inverting Eq.~\maineqref{eq:working}{2} within each mark count, and replaying the complete policy reproduces the saved inclusion probabilities. All four policies in Tab.~\ref{tab:ablation} share the frozen leverage $\bar h_i$, the square-root rule of Eq.~\maineqref{eq:allocation}{3}, the floor $\pi_{\min}b=0.0025$, the cap \ProbCap{} and the expected budget of 3,950.15 candidates; no label, auxiliary, head or detector is refitted. The image-uncertainty control sets $u_i=p_i$ and $v_i=\max\{p_i(1-p_i),\varepsilon_v\}$, and the marks-only control fixes $p_i=\MarksOnlyP$, giving $u_i=0.175$, 0.725 and 0.95 for zero, one and two marks. Both controls were defined before their reach was computed, and neither is optimized. Reach is exact (Eq.~\maineqref{eq:reach}{4}), so no audit draws are simulated. Intervals resample the \NFit{} training patients \NBoot{} times, sum the reach within patients and divide by the resampled number of omissions. They condition on the fixed scores and allocation probabilities, omit refitting, calibration and design variation, and are not adjusted for multiplicity. Since $p_i$ is an out-of-fold propensity of either weak mark rather than a calibrated probability of a reference target, the image-uncertainty control also changes the probability mapping, and its difference from the complete policy is a practical ablation rather than an identified effect of the marks.

\begin{table}[t]
\centering\small
\setlength{\tabcolsep}{2.2pt}
\caption{Sources of the acquisition gain at half budget (post hoc replay, no retraining). Top: expected training targets reached among the \NOmitTrain{} omissions, the \NOmitUnmarked{} omissions without any weak mark and all \NTargetTrain{} targets. Bottom: complete policy minus each control, in expected omissions and in percentage points of omission reach, with conditional patient-bootstrap \CILevel\% intervals.}
\label{tab:ablation}
\begin{tabular}{lrrr}
\toprule
Policy & Omit.\ (92) & No mark (35) & All (541) \\
\midrule
Image tail, constant $v_i$ & 64.2 & 22.6 & 429.2 \\
Marks only, $p_i=0.5$ & 69.7 & 23.5 & 341.8 \\
Image uncertainty, $u_i=p_i$ & 76.1 & 27.8 & 479.1 \\
Paired tail (complete) & 79.4 & 27.4 & 484.0 \\
\bottomrule
\end{tabular}

\end{table}

\paragraph{Leverage terms.}
This post hoc diagnostic holds the fitted model, the paired proxy variance, the 378,371-candidate roster, the probability bounds and the expected budget of 3,950.15 candidate decisions fixed, and varies only the leverage factor of Eq.~\maineqref{eq:leverage}{1}. It covers all \NFit{} training patients, \NTargetTrain{} targets and \NOmitTrain{} omissions, of which \NOmitUnmarked{} carry no weak mark; it was specified before the contrasts were computed, and replaying the full leverage reproduces the original allocation exactly. Rank leverage is $e^{-(r_i-1)/\tau}$, threshold leverage is $\sigma((s_i-t_0+\delta)/\omega)$ with the saved $t_0=1.566$, and full leverage takes their maximum. Every factor is floored at $\varepsilon_h$ and multiplied by $\sqrt{\max\{u_i(1-u_i),\varepsilon_v\}}$, and each variant solves its own multiplier to attain the same expected budget, with probabilities clipped to $[0.0025, \ProbCap]$. Constant leverage keeps the same mark-conditioned variance, so it is neither uniform sampling nor the image-uncertainty control above. In omission reach, full leverage exceeds constant leverage by 79.5 points [74.0, 84.1], the threshold term alone by 15.4 points [9.6, 21.3] and the rank term alone by only $+0.047$ points [$-0.261$, 0.593] (Tab.~\ref{tab:leverage}). The intervals use \NBoot{} paired patient-bootstrap resamples conditional on the fitted scores, threshold and allocation; they are pointwise and unadjusted, and describe reach on known training references, not generalization or detector accuracy. All four variants are reported. Full leverage reaches slightly fewer omissions without any weak mark than the rank term alone (27.39 \vs 27.46); its higher all-target reach (483.98 \vs 476.83) is a distinct endpoint and does not establish an advantage in omission reach. Both variants open nearly every training patient, so the comparison concerns candidate decisions, not proportional savings in reading time.

\begin{table}[t]
\centering\small
\setlength{\tabcolsep}{2.6pt}
\caption{Leverage factorial at half budget: expected training targets reached among the \NOmitTrain{} omissions, the \NOmitUnmarked{} omissions without any weak mark and all \NTargetTrain{} targets, and the expected number of patients with at least one audit.}
\label{tab:leverage}
\begin{tabular}{lrrrr}
\toprule
Leverage & Omit.\ (92) & No mark (35) & All (541) & Patients \\
\midrule
Constant & 6.27 & 1.79 & 23.76 & 443.04 \\
Rank only & 79.33 & 27.46 & 476.83 & 444.98 \\
Threshold only & 65.22 & 21.62 & 441.39 & 442.44 \\
Full (maximum) & 79.38 & 27.39 & 483.98 & 444.97 \\
\bottomrule
\end{tabular}

\end{table}

\FloatBarrier
\section{Calibration Reference Analysis}
\label{sec:s-ref}

\paragraph{Alert-rate bound.}
Let $S_{n+1}$ be the maximum of a new negative patient, and suppose that $S_1,\dots,S_{n+1}$ are exchangeable and have no ties. The rank of $S_{n+1}$ among them is then uniform on $\{1,\dots,n+1\}$, and $S_{n+1}>S_{(k)}$ holds exactly when this rank exceeds $k$. Hence $\Pr(S_{n+1}>T_\alpha)=(n+1-k)/(n+1)\le\alpha$ for $k=\lceil(n+1)(1-\alpha)\rceil$, and ties can only lower this probability. The bound is marginal over the calibration sample and the new patient, and it applies to the population of which the calibration negatives are exchangeable draws. A pool $P=U\cup E$ calibrates for the mixture of patients with and without an official consensus nodule, and this choice of population is what the score-fixed replay of Sec.~\mainref{sec:tail}{5} varies.

\paragraph{Negative pools.}
The C/D pool $P$ contains the patients without a concordant C/D contour under our geometry, and the augmented pool $U$ additionally requires that the patient has no entry in the official LUNA16 annotation file. For visualization, candidates are matched to official nodules with the same tolerance as for proposals, namely the greater of \MinRadiusMM\,mm and half the official diameter. Patient eligibility, by contrast, depends only on the presence of an official nodule in the scan, whether or not a candidate matches it. Since changing the negative pool leaves the localization targets unchanged, thresholds can be compared without altering the recall denominator, and under both thresholds the alert denominator consists of the \UTest{} test patients that are negative under both definitions.

\begin{table*}[t]
\centering\small
\caption{Reference effects by feature family. Each row summarizes distinct fixed score vectors. Score vectors: distinct score vectors among the \NHeadRows{} fitted configurations; Lower threshold: score vectors whose threshold decreases; Past size band: displacements beyond the central \CILevel\% of random removals of the same size; More omissions: score vectors with strictly more omission detections under the lower threshold; Median ranks: median threshold displacement.}
\label{tab:families}
\begin{tabular}{lrrrrr}
\toprule
Features & Score vectors & Lower threshold & Past size band & More omissions & Median ranks \\
\midrule
Descriptors & 24 & 24 & 4 & 11 & 3 \\
Multi-view & 2 & 2 & 2 & 2 & 6.5 \\
3D & 257 & 254 & 196 & 238 & 6 \\
\bottomrule
\end{tabular}

\end{table*}

\begin{table*}[t]
\centering\small
\caption{Paired-tail differences under each calibration reference. The six matched fits are averaged within each patient. C/D: pool $P$; + official: pool $U$. Both pools are evaluated on the same original targets and the \UTest{} common negative test patients. Differences and patient-bootstrap intervals are in percentage points.}
\label{tab:alt}
\begin{tabular}{llrrr}
\toprule
Baseline & Cal. pool & Omit. diff. [95\%] & All diff. [95\%] & Alert diff. [95\%] \\
\midrule
Mark IPW & C/D & -4.0 [-7.0, -1.4] & +3.0 [+1.3, +4.9] & -1.5 [-3.1, +0.0] \\
Mark IPW & + official & +3.6 [-1.4, +9.5] & +6.0 [+3.8, +8.4] & +1.1 [-0.8, +3.4] \\
Mark, same head & C/D & -4.4 [-8.3, -0.8] & -1.3 [-3.0, +0.4] & -2.0 [-4.2, -0.4] \\
Mark, same head & + official & +4.4 [-0.9, +10.0] & +4.8 [+3.0, +6.7] & +0.9 [-0.7, +2.7] \\
First read & C/D & -6.3 [-15.1, +0.0] & -6.2 [-9.3, -3.4] & +0.7 [+0.2, +1.5] \\
First read & + official & -6.0 [-12.2, -1.6] & -3.8 [-6.2, -1.4] & +1.7 [-0.6, +4.3] \\
\bottomrule
\end{tabular}

\end{table*}

\paragraph{Coverage of fitted heads and an additional architecture.}
For every locked fit, we first reproduce the original thresholds and all C/D counts at $\alpha=\AlphaNom$ and then evaluate the augmented pool (Tab.~\ref{tab:families}). To test reference sensitivity in a second learned representation, we additionally use the epoch-16 checkpoint of the 2.5D model from the first development screen. Its input consists of twelve planes centred on the candidate at two physical fields of view, processed by two blocks each of 24, 40 and 64 channels with two intermediate pooling operations; the embedding concatenates the central and global means before a 32-channel projection, and at inference it is averaged over the patch and its in-plane flips. With this checkpoint frozen, we fit first-read and full-reference linear heads on the original fitting set with the common ridge penalty. These fits use the same patients and candidate source as the main experiments, and identical score vectors are merged when counting distinct detectors. Table~\ref{tab:alt} gives the paired-tail differences from each baseline under both pools.

\paragraph{Size-matched random removal.}
Within each of the two calibration subsets, we draw a random subset of negative patients of the same size as the augmented pool in that subset, and repeat this \NRandom{} times with common random draws across detectors. For each draw, we take the \URank{}th order statistic, as for the augmented pool, and locate the resulting score in the sorted original pool, breaking ties deterministically by roster order. The resulting distribution describes how far the threshold would move if patients were removed uniformly at random within each subset at the same sample size, and the observed displacement is expressed on the same rank scale. As a descriptive tail area, we report the fraction of random displacements that are at least as large as the observed one, with a one-draw correction.

\paragraph{Uncertainty.}
Paired test-patient intervals use \NBoot{} bootstrap samples after averaging the matched audit fits within each patient. When the two references are compared, the scores, calibration pools and thresholds are held fixed during this bootstrap, so the intervals describe variation across test patients under those settings, whereas the random-removal distribution varies only the calibration membership. The reference analysis was specified after the locked comparison had been scored, whereas the training and allocation protocol defines the locked comparison itself.

\paragraph{Alternative reader pairs.}
We re-parse the four LIDC sessions with the same geometry and, for each of the six unordered reader pairs, reconstruct the pair consensus without using any score, after verifying that the C/D pair reproduces the locked target count on every retained scan. Each pair's negatives are then used for calibration, and the resulting thresholds are applied to the original targets (Tab.~\ref{tab:roles}). Test alerts are counted on the \RoleTest{} patients that are negative under all six pairs and under the official list. This analysis isolates the sensitivity to the reader assignment, although every pair reuses the same unblinded records.

\begin{table}[t]
\centering\small
\caption{Calibration on each reader pair at the anchor draw. Cal.\ neg.: negative calibration patients under the pair; cells give omission detections out of \NOmitTest{} / alerts among the \RoleTest{} common negative test patients.}
\label{tab:roles}
\begin{tabular}{lrrrr}
\toprule
Pair & Cal. neg. & First read & Mark IPW & Paired tail \\
\midrule
AB & 66 & 6 / 0 & 3 / 1 & 1 / 0 \\
AC & 72 & 6 / 0 & 3 / 1 & 1 / 0 \\
AD & 63 & 6 / 0 & 3 / 1 & 5 / 1 \\
BC & 68 & 6 / 3 & 4 / 1 & 2 / 0 \\
BD & 65 & 6 / 0 & 3 / 1 & 2 / 0 \\
CD & 67 & 4 / 0 & 3 / 1 & 1 / 0 \\
\bottomrule
\end{tabular}

\end{table}

\paragraph{Masking disputed candidates.}
Both maxima that trigger an alert on a C/D-negative test patient are disputed candidates, although only \NSingleContour{} of the \NCalNegCand{} candidates on C/D-negative calibration scans are disputed (main paper, Sec.~\mainref{sec:misses}{4.5}). In a locked sensitivity analysis of the selected method at the anchor draw, we removed the disputed candidates from the roster. At the original threshold, this removes both negative alerts and keeps the one recovered omission. Recalibrating after the removal gives 12 of \NOmitTest{} recovered omissions and five alerts among the \PTest{} C/D-negative test patients. Unlike the score-fixed replay, this intervention changes which candidates are eligible, so it cannot separate the effect of patient membership from that of the roster; the replay of Sec.~\mainref{sec:tail}{5} keeps the roster and every score fixed.

\FloatBarrier
\section{Simulation Study of the Estimator}
\label{sec:s-sim}

In a controlled simulation, we fix a finite synthetic population and its auxiliary fit and vary the audit draws and the dependence between the two weak readers. The estimand is the proportion of positives among candidates left unmarked by the first reader. At the base budget, the augmented estimate of this proportion reduces the root-mean-square error by about 10--12\% relative to the Horvitz--Thompson estimate. Note that the localization experiments in the main paper and in the bookmark-cap setting target different estimands.

\FloatBarrier
